% GENERATED FILE. Edit canonical lessons, then run npm run generate:books.
% canonical-source-hash: fnv1a64:c34d8255056be9dc
% canonical-lessons: SA-C101-kankatika, SA-C101-phenakam, SA-C101-pronchanam, SA-C101-upanetram, SA-C101-syutah

\chapter{Comb, Soap, Towel, Glasses, Bag}
\label{ch:sa-comb-soap-towel-glasses-bag}
\hlchaptermodality{101}

\begin{chapteropening}
\textbf{By the end of this chapter:} \emph{I can say a comb, soap, a towel, glasses and a bag.}

Complete the last lesson of chapter 101: I can say a comb, soap, a towel, glasses and a bag.
\end{chapteropening}

\section[\texorpdfstring{\sk{कंकतिका} (kaṅkatikā)}{कंकतिका (kaṅkatikā)}]{\sk{कंकतिका} (kaṅkatikā) — a comb}

\label{lesson:SA-C101-kankatika}



\begin{quote}
Before the new one: say the Sanskrit for a needle, then the Sanskrit for thread.
\end{quote}

\subsection*{You'll want to know: \sk{कंकतिका}}
\textbf{\sk{कंकतिका}} — \emph{kaṅkatikā} — \textquotedblleft{}a comb\textquotedblright{}.

\begin{grammarlens}[title={What you've built}]
One more word for this chapter.
\end{grammarlens}

\subsection*{Your turn}
\begin{itemize}
\raggedright
  \item \emph{Say it:} \emph{kaṅkatikā}
  \item \emph{Say it:} \emph{kaṅkatikā}, once more
  \item \emph{Say it:} say \emph{kaṅkatikā}
  \item \emph{Recall:} say the Sanskrit for a needle, then the Sanskrit for thread, then say \emph{kaṅkatikā} again
\end{itemize}

\begin{tcolorbox}[breakable,colback=teal!4,colframe=teal!35!black,title={Before you move on}]
What does \sk{कंकतिका} mean? (\textquotedblleft{}A comb\textquotedblright{}.) Say it once more.
\end{tcolorbox}
\section[\texorpdfstring{\sk{फेनकम्} (phenakam)}{फेनकम् (phenakam)}]{\sk{फेनकम्} (phenakam) — soap}

\label{lesson:SA-C101-phenakam}



\begin{quote}
Before the new one: say the Sanskrit for thread, then the Sanskrit for a comb.
\end{quote}

\subsection*{You'll want to know: \sk{फेनकम्}}
\textbf{\sk{फेनकम्}} — \emph{phenakam} — \textquotedblleft{}soap\textquotedblright{}.

\begin{grammarlens}[title={What you've built}]
One more word for this chapter.
\end{grammarlens}

\subsection*{Your turn}
\begin{itemize}
\raggedright
  \item \emph{Say it:} \emph{phenakam}
  \item \emph{Say it:} \emph{phenakam}, once more
  \item \emph{Say it:} say \emph{phenakam}
  \item \emph{Recall:} say the Sanskrit for thread, then the Sanskrit for a comb, then say \emph{phenakam} again
\end{itemize}

\begin{tcolorbox}[breakable,colback=teal!4,colframe=teal!35!black,title={Before you move on}]
What does \sk{फेनकम्} mean? (\textquotedblleft{}Soap\textquotedblright{}.) Say it once more.
\end{tcolorbox}
\section[\texorpdfstring{\sk{प्रोञ्छनम्} (proñchanam)}{प्रोञ्छनम् (proñchanam)}]{\sk{प्रोञ्छनम्} (proñchanam) — a towel}

\label{lesson:SA-C101-pronchanam}



\begin{quote}
Before the new one: say the Sanskrit for a comb, then the Sanskrit for soap.
\end{quote}

\subsection*{You'll want to know: \sk{प्रोञ्छनम्}}
\textbf{\sk{प्रोञ्छनम्}} — \emph{proñchanam} — \textquotedblleft{}a towel\textquotedblright{}.

\begin{grammarlens}[title={What you've built}]
One more word for this chapter.
\end{grammarlens}

\subsection*{Your turn}
\begin{itemize}
\raggedright
  \item \emph{Say it:} \emph{proñchanam}
  \item \emph{Say it:} \emph{proñchanam}, once more
  \item \emph{Say it:} say \emph{proñchanam}
  \item \emph{Recall:} say the Sanskrit for a comb, then the Sanskrit for soap, then say \emph{proñchanam} again
\end{itemize}

\begin{tcolorbox}[breakable,colback=teal!4,colframe=teal!35!black,title={Before you move on}]
What does \sk{प्रोञ्छनम्} mean? (\textquotedblleft{}A towel\textquotedblright{}.) Say it once more.
\end{tcolorbox}
\section[\texorpdfstring{\sk{उपनेत्रम्} (upanetram)}{उपनेत्रम् (upanetram)}]{\sk{उपनेत्रम्} (upanetram) — glasses, spectacles}

\label{lesson:SA-C101-upanetram}



\begin{quote}
Before the new one: say the Sanskrit for soap, then the Sanskrit for a towel.
\end{quote}

\subsection*{You'll want to know: \sk{उपनेत्रम्}}
\textbf{\sk{उपनेत्रम्}} — \emph{upanetram} — \textquotedblleft{}glasses, spectacles\textquotedblright{}.

\begin{grammarlens}[title={What you've built}]
One more word for this chapter.
\end{grammarlens}

\subsection*{Your turn}
\begin{itemize}
\raggedright
  \item \emph{Say it:} \emph{upanetram}
  \item \emph{Say it:} \emph{upanetram}, once more
  \item \emph{Say it:} say \emph{upanetram}
  \item \emph{Recall:} say the Sanskrit for soap, then the Sanskrit for a towel, then say \emph{upanetram} again
\end{itemize}

\begin{tcolorbox}[breakable,colback=teal!4,colframe=teal!35!black,title={Before you move on}]
What does \sk{उपनेत्रम्} mean? (\textquotedblleft{}Glasses, spectacles\textquotedblright{}.) Say it once more.
\end{tcolorbox}
\section[\texorpdfstring{\sk{स्यूतः} (syūtaḥ)}{स्यूतः (syūtaḥ)}]{\sk{स्यूतः} (syūtaḥ) — a bag}

\label{lesson:SA-C101-syutah}



\begin{quote}
Before the new one: say the Sanskrit for a towel, then the Sanskrit for glasses.
\end{quote}

\subsection*{You'll want to know: \sk{स्यूतः}}
\textbf{\sk{स्यूतः}} — \emph{syūtaḥ} — \textquotedblleft{}a bag\textquotedblright{}.

\begin{grammarlens}[title={What you've built}]
One more word for this chapter.
\end{grammarlens}

\subsection*{Your turn}
\begin{itemize}
\raggedright
  \item \emph{Say it:} \emph{syūtaḥ}
  \item \emph{Say it:} \emph{syūtaḥ}, once more
  \item \emph{Say it:} say \emph{syūtaḥ}
  \item \emph{Recall:} say the Sanskrit for a towel, then the Sanskrit for glasses, then say \emph{syūtaḥ} again
\end{itemize}

\begin{tcolorbox}[breakable,colback=teal!4,colframe=teal!35!black,title={Before you move on}]
What does \sk{स्यूतः} mean? (\textquotedblleft{}A bag\textquotedblright{}.) Say it once more.
\end{tcolorbox}
